% Nature Journal Article Template
% Drafting aid for Nature; other Nature Portfolio journals differ
% Reviewed: 2026-10-01
% Generic writing scaffold, not an official Nature template.
% Check the exact journal, article type, and submission stage at:
% https://www.nature.com/nature/for-authors/initial-submission

\documentclass[12pt]{article}

% Packages
\usepackage[margin=2.5cm]{geometry}
\usepackage{times}
\usepackage{graphicx}
\usepackage{amsmath}
\usepackage{amssymb}
\usepackage{hyperref}
\usepackage{lineno}  % Line numbers for review
\usepackage[super]{natbib}  % Superscript citations

% Line numbering (required for submission)
\linenumbers

% Title and Authors
\title{Insert Your Title Here: Concise and Descriptive}

\author{
First Author\textsuperscript{1}, Second Author\textsuperscript{1,2}, Third Author\textsuperscript{2,*}
}

\date{}

\begin{document}

\maketitle

% Affiliations
\noindent
\textsuperscript{1}Department Name, Institution Name, City, State/Province, Postal Code, Country \\
\textsuperscript{2}Second Department/Institution \\
\textsuperscript{*}Correspondence: [email protected]

% Abstract
\begin{abstract}
\noindent
Write a referenced summary paragraph, ideally no more than 200 words, for a Nature Article. Summarize the actual findings, significance, and conclusions. The abstract should be self-contained and understandable without reading the full paper. Focus on what you did, what you found, and why it matters. Avoid jargon and abbreviations where possible.
\end{abstract}

% Main Text
\section*{Introduction}
% 2-3 paragraphs setting the context
Provide background on the research area, establish the importance of the problem, and identify the knowledge gap your work addresses. Nature papers should emphasize broad significance beyond a narrow specialty.

State your main research question or objective clearly.

Briefly preview your approach and key findings.

\section*{Results}
% Primary results section
% Organize by finding, not by experiment
% Reference figures/tables as you describe results

\subsection*{First major finding}
Describe your first key result. Reference Figure~\ref{fig:example} to support your findings.

\begin{figure}[ht]
\centering
% Include your figure here
% \includegraphics[width=0.7\textwidth]{figure1.pdf}
\caption{{\bf Figure title in bold.} Detailed figure caption explaining what is shown, experimental conditions, sample sizes (n), statistical tests, and significance levels. Panels should be labeled (a), (b), etc. if multiple panels are present.}
\label{fig:example}
\end{figure}

\subsection*{Second major finding}
Describe your second key result objectively, without interpretation.

\subsection*{Third major finding}
Describe additional results as needed.

\section*{Discussion}
% Interpret results, compare to literature, acknowledge limitations

\subsection*{Main findings and interpretation}
Summarize your key findings and explain their significance. How do they advance our understanding?

\subsection*{Comparison to previous work}
Compare and contrast your results with existing literature\cite{example2023}.

\subsection*{Implications}
Discuss the broader implications of your work for the field and beyond.

\subsection*{Limitations and future directions}
Honestly acknowledge limitations and suggest future research directions.

\section*{Conclusions}
Provide a concise conclusion summarizing the main take-home messages of your work.

\section*{Data availability}
State the actual repository/accession, supporting files, or justified restrictions. Follow required repository deposition rules; available-on-request is not a universal alternative.

\section*{Code availability}
If applicable, provide information on code availability (GitHub, Zenodo, etc.).

\section*{Acknowledgements}
Acknowledge non-author contributions accurately.

\section*{Funding}
State the actual funding sources and grant identifiers in the required separate funding statement.

\section*{Author contributions}
Describe the actual contributions of each author, for example study design, methods, investigation, analysis, and writing.

\section*{Competing interests}
Declare any financial or non-financial competing interests. If none, state "The authors declare no competing interests."

% References
% Optionally use a verified bibliography style and your real .bib file instead
% of the manual placeholder list below; do not enable both approaches.
% \bibliographystyle{naturemag}
% \bibliography{references}

% Alternatively, manually format references:
\begin{thebibliography}{99}

\bibitem{example2023}
Replace with the verified bibliographic record supporting the cited claim.

% Add more references as needed

\end{thebibliography}

% Figure Legends (if not included with figures)
\section*{Figure Legends}

\textbf{Figure 1 | Figure title.} Comprehensive figure legend describing all panels, experimental conditions, sample sizes, and statistical analyses.

\textbf{Figure 2 | Second figure title.} Another detailed legend.

\section*{Methods}
% Detailed methods allowing reproducibility
% Nature house organization places Methods after figure legends in the main file.
% Confirm the current submission-stage organization.

\subsection*{Experimental design}
Describe overall experimental design, including controls.

\subsection*{Sample preparation}
Detail procedures for sample collection, preparation, and handling.

\subsection*{Data collection}
Describe instrumentation, measurement procedures, and data collection protocols.

\subsection*{Data analysis}
Explain analytical methods, statistical tests, and software used. State sample sizes, replication, and significance thresholds.

\subsection*{Ethical approval}
Include relevant ethical approval statements (human subjects, animal use, biosafety).

% Extended Data Figures (optional)
\section*{Extended Data}

\textbf{Extended Data Figure 1 | Supplementary data title.} Description of supplementary figure supporting main findings.

\end{document}

% Notes for Authors:
% 1. Exact length depends on article type and stage; published-layout examples are not a universal word cap
% 2. Use superscript numbered citations (1, 2, 3)
% 3. Use the current figure guide at final physical size; retain editable artwork
% 4. Nature encourages combined text/figures for initial review; production differs
% 5. Follow the current stage-specific spacing instructions
% 6. Include line numbers using \linenumbers
% 7. Follow Nature's specific author guidelines for your target journal
% 8. Methods section can be quite detailed and placed after main text
% 9. Check word limits and specific requirements for your Nature family journal

